% =====================================================================
% Ablation figures — Stage 2 results (val vs test overlays)
% Generated for exp8_ablation_model_size
%
% Required preamble (add once near the top of your main .tex):
%   \usepackage{graphicx}
%   \usepackage{subcaption}
%
% Folder layout assumed at the project root (as uploaded):
%   nn/        qrf/        xgb/        xgb_1sub/
%
% If you uploaded everything inside a parent folder (e.g. results/ or figs/),
% add this ONE line above the figures and edit the path:
%   \graphicspath{{results/}}
% =====================================================================


% ===================== NN ablation =====================
\begin{figure}[t]
  \centering
  \begin{subfigure}[t]{0.48\textwidth}\centering
    \includegraphics[width=\linewidth]{nn/nn_hidden_sizes_w_sum_rmse_2026_06_22_12_05_40.png}
    \caption{Sum of RMSEs}
    \label{fig:abl-nn-sr}
  \end{subfigure}\hfill
  \begin{subfigure}[t]{0.48\textwidth}\centering
    \includegraphics[width=\linewidth]{nn/nn_hidden_sizes_w_sigma_opt_2026_06_22_12_05_40.png}
    \caption{$\sigma_{\mathrm{opt}}$}
    \label{fig:abl-nn-so}
  \end{subfigure}
  \caption{NN ablation over hidden width $h \in \{16, 32, 64, 128, 256, 512, 1024\}$,
  averaged over $D = 10$ data seeds and $I = 10$ init seeds. Sum of RMSEs
  falls monotonically, but parameter uncertainty $\sigma_{\mathrm{opt}}$ reaches
  its minimum at $h \in [64, 128]$ and rises sharply for wider networks.}
  \label{fig:abl-nn}
\end{figure}


% ===================== QRF ablation: number of trees =====================
\begin{figure}[t]
  \centering
  \begin{subfigure}[t]{0.48\textwidth}\centering
    \includegraphics[width=\linewidth]{qrf/qrf_n_estimators_sum_rmse_2026_06_22_12_05_40.png}
    \caption{Sum of RMSEs}
  \end{subfigure}\hfill
  \begin{subfigure}[t]{0.48\textwidth}\centering
    \includegraphics[width=\linewidth]{qrf/qrf_n_estimators_sigma_opt_2026_06_22_12_05_40.png}
    \caption{$\sigma_{\mathrm{opt}}$}
  \end{subfigure}
  \caption{QRF ablation over number of trees,
  $n_{\mathrm{est}} \in \{50, 100, 150, 200, 300, 500\}$,
  with the other QRF hyperparameters fixed to their stage-1 best values.}
  \label{fig:abl-qrf-nest}
\end{figure}


% ===================== QRF ablation: min samples per leaf =====================
\begin{figure}[t]
  \centering
  \begin{subfigure}[t]{0.48\textwidth}\centering
    \includegraphics[width=\linewidth]{qrf/qrf_min_samples_leaf_sum_rmse_2026_06_22_12_05_40.png}
    \caption{Sum of RMSEs}
  \end{subfigure}\hfill
  \begin{subfigure}[t]{0.48\textwidth}\centering
    \includegraphics[width=\linewidth]{qrf/qrf_min_samples_leaf_sigma_opt_2026_06_22_12_05_40.png}
    \caption{$\sigma_{\mathrm{opt}}$}
  \end{subfigure}
  \caption{QRF ablation over the minimum number of samples per leaf,
  $\mathrm{min\_samples\_leaf} \in \{2, 5, 10, 20, 50\}$, with the other QRF
  hyperparameters fixed to their stage-1 best values.}
  \label{fig:abl-qrf-msl}
\end{figure}


% ===================== XGB ablation: number of estimators =====================
\begin{figure}[t]
  \centering
  \begin{subfigure}[t]{0.48\textwidth}\centering
    \includegraphics[width=\linewidth]{xgb/xgb_n_estimators_sum_rmse_2026_06_22_12_05_40.png}
    \caption{Sum of RMSEs}
  \end{subfigure}\hfill
  \begin{subfigure}[t]{0.48\textwidth}\centering
    \includegraphics[width=\linewidth]{xgb/xgb_n_estimators_sigma_opt_2026_06_22_12_05_40.png}
    \caption{$\sigma_{\mathrm{opt}}$}
  \end{subfigure}
  \caption{XGB ablation over number of estimators,
  $n_{\mathrm{est}} \in \{50, 100, 200, 300, 500\}$, with $\mathrm{max\_depth}$,
  $\mathrm{learning\_rate}$ and $\mathrm{subsample}$ fixed to their stage-1
  best values.}
  \label{fig:abl-xgb-nest}
\end{figure}


% ===================== XGB ablation: max depth =====================
\begin{figure}[t]
  \centering
  \begin{subfigure}[t]{0.48\textwidth}\centering
    \includegraphics[width=\linewidth]{xgb/xgb_max_depth_sum_rmse_2026_06_22_12_05_40.png}
    \caption{Sum of RMSEs}
  \end{subfigure}\hfill
  \begin{subfigure}[t]{0.48\textwidth}\centering
    \includegraphics[width=\linewidth]{xgb/xgb_max_depth_sigma_opt_2026_06_22_12_05_40.png}
    \caption{$\sigma_{\mathrm{opt}}$}
  \end{subfigure}
  \caption{XGB ablation over $\mathrm{max\_depth} \in \{1, 2, 3, 4, 5, 7, 10\}$,
  with $n_{\mathrm{est}}$, $\mathrm{learning\_rate}$ and $\mathrm{subsample}$
  fixed to their stage-1 best values.}
  \label{fig:abl-xgb-md}
\end{figure}


% ===================== XGB-1sub ablation: number of estimators =====================
\begin{figure}[t]
  \centering
  \begin{subfigure}[t]{0.48\textwidth}\centering
    \includegraphics[width=\linewidth]{xgb_1sub/xgb-1sub_n_estimators_sum_rmse_2026_06_22_12_05_40.png}
    \caption{Sum of RMSEs}
  \end{subfigure}\hfill
  \begin{subfigure}[t]{0.48\textwidth}\centering
    \includegraphics[width=\linewidth]{xgb_1sub/xgb-1sub_n_estimators_sigma_opt_2026_06_22_12_05_40.png}
    \caption{$\sigma_{\mathrm{opt}}$}
  \end{subfigure}
  \caption{XGB-1sub ablation over number of estimators, with
  $\mathrm{subsample} = 1$ (deterministic boosting). Here
  $\sigma_{\mathrm{opt}} = 0$ by construction; the curve confirms that
  removing stochastic subsampling removes parameter uncertainty entirely.}
  \label{fig:abl-xgb1sub-nest}
\end{figure}


% ===================== XGB-1sub ablation: max depth =====================
\begin{figure}[t]
  \centering
  \begin{subfigure}[t]{0.48\textwidth}\centering
    \includegraphics[width=\linewidth]{xgb_1sub/xgb-1sub_max_depth_sum_rmse_2026_06_22_12_05_40.png}
    \caption{Sum of RMSEs}
  \end{subfigure}\hfill
  \begin{subfigure}[t]{0.48\textwidth}\centering
    \includegraphics[width=\linewidth]{xgb_1sub/xgb-1sub_max_depth_sigma_opt_2026_06_22_12_05_40.png}
    \caption{$\sigma_{\mathrm{opt}}$}
  \end{subfigure}
  \caption{XGB-1sub ablation over $\mathrm{max\_depth} \in \{1, 2, 3, 4, 5, 7, 10\}$,
  with $\mathrm{subsample} = 1$ (deterministic boosting). $\sigma_{\mathrm{opt}} = 0$
  by construction at every depth.}
  \label{fig:abl-xgb1sub-md}
\end{figure}
